%% =====================================================================
%%  B6-T-19-05 -- what B6 contributes to "The Conscience Attack"
%%  -------------------------------------------------------------------
%%  LAYER A: APPEND-ONLY. Written once at the entry's write, then left
%%  alone. If the source has more to say later, add a bullet -- do not
%%  rewrite. Material found ONLY here goes in the `unique` slot with
%%  @unique set to yes, which makes it undeletable (rule R10).
%% =====================================================================
%% @kind: source
%% @id: B6-T-19-05
%% @book: B6
%% @topic: T-19-05
%% @locator: ch. 9 (guilt-tripping and shaming), pp. 116--120; ch. 7 (Janice and Bill), pp. 92--105
%% @status: distilled
%% @unique: no
%% @integrated: yes
%% @updated: 2026-09-19

\dossier{B6}{T-19-05}{\bookshort{B6}}{ch. 9 (guilt-tripping and shaming), pp. 116--120; ch. 7 (Janice and Bill), pp. 92--105}

\begin{distilled}
  \item Manipulators use what they know to be the greater conscientiousness of their victims to keep them self-doubting, anxious and submissive: the more conscientious the potential victim, the more effective guilt is as a weapon.
  \item Shaming works beside guilt: subtle sarcasm and put-downs increase fear and self-doubt, fostering a standing sense of inadequacy in the weaker party and thereby maintaining dominance.
  \item The Janice and Bill case: relapses timed to her exit thinking convert her independence into guilt; the staged overdose theatre completes the training after her anger has been taught to doubt its own legitimacy.
  \item The caretaker profile is the target profile: she hates thinking of herself as the bad guy, is riddled with guilt when she feels selfish, and therefore pays the attack's full price.
\end{distilled}
